% !TEX root = ../meu-trabalho.tex
% -----------------------------------------------------------------------------
% Estudos relacionados
% -----------------------------------------------------------------------------

\chapter{Estudos relacionados}
\label{chap_estudos_relacionados}

\section{Estudos sobre machine learning aplicado à otimização de portfólios}
\label{sec_estudos_sobre_machine_learning_aplicado_a_otimizacao_de_portfolios}

Os estudos analisados sobre otimização de portfólios mostram que modelos preditivos vêm sendo combinados com métodos tradicionais de alocação de recursos. Em geral, essa combinação ocorre em duas etapas: primeiro, algoritmos de machine learning ou deep learning são usados para prever preços, retornos ou movimentos futuros dos ativos; depois, essas previsões são incorporadas a modelos de seleção e otimização de carteiras, especialmente ao modelo média-variância de Markowitz \cite{chen2021mean,ma2021portfolio,wang2020portfolio}. Embora esses estudos tratem principalmente de ações e índices, sua lógica metodológica se aproxima do tema deste anteprojeto, pois combina previsão algorítmica e diversificação de portfólios.

A previsão aparece como uma etapa importante porque fornece informações para a montagem posterior da carteira. \citeonline{chen2021mean} combinam previsão de preços por machine learning com a definição dos pesos da carteira pelo modelo média-variância. \citeonline{ma2021portfolio} também relacionam previsões de retorno à otimização de portfólios, destacando que modelos de machine learning e deep learning podem apresentar desempenho preditivo mais favorável do que modelos tradicionais de séries temporais. De forma semelhante, \citeonline{wang2020portfolio} utilizam LSTM para capturar dependências de longo prazo nos retornos dos ativos e, em seguida, formar carteiras otimizadas pelo critério média-variância.

Além da previsão, os estudos destacam a pré-seleção de ativos antes da alocação final. \citeonline{chen2021mean} selecionam ativos com maior potencial de retorno a partir dos preços previstos, enquanto \citeonline{ma2021portfolio} utilizam modelos preditivos para filtrar ativos antes de aplicar os modelos média-variância e ômega. \citeonline{wang2020portfolio} também enfatizam a seleção inicial dos ativos que irão compor a carteira. Assim, a pré-seleção por algoritmos funciona como uma etapa intermediária entre a previsão e a otimização, direcionando a alocação para ativos com maior potencial estimado \cite{chen2021mean,ma2021portfolio,wang2020portfolio}.

Quanto ao desempenho, os estudos sugerem que a combinação entre previsão e otimização pode gerar resultados mais favoráveis do que estratégias sem etapa preditiva. \citeonline{chen2021mean}, em estudo aplicado à Bolsa de Xangai, reportam melhores resultados em retorno e risco quando comparados a métodos tradicionais sem previsão e a diferentes benchmarks. \citeonline{ma2021portfolio} identificam desempenho mais favorável nos modelos média-variância e ômega quando combinados com previsões de Random Forest. \citeonline{wang2020portfolio}, por sua vez, indicam que a estratégia LSTM + MV apresentou melhorias em retorno cumulativo e retorno ajustado ao risco na maior parte dos períodos analisados.

\citeonline{padhi2022intelligent} também tratam dessa integração ao combinar seleção de portfólio e aprendizado de máquina. Nesse estudo, a abordagem média-variância é usada para reduzir o risco por meio da construção de uma carteira ótima, enquanto técnicas de aprendizado online são aplicadas para prever movimentos futuros de preços. A partir do material disponível, \citeonline{paiva2019decision} também podem ser associados a esse eixo, pois tratam de uma abordagem de fusão entre machine learning e seleção de portfólio para tomada de decisão em negociações financeiras. Contudo, os dados fornecidos não apresentam detalhes suficientes sobre seus métodos e resultados para uma comparação mais aprofundada.

De modo geral, os estudos analisados indicam que o machine learning tem sido usado na otimização de portfólios principalmente para prever retornos ou preços, selecionar ativos promissores e buscar melhor desempenho após a alocação matemática \cite{chen2021mean,ma2021portfolio,wang2020portfolio,padhi2022intelligent,paiva2019decision}. Ainda assim, os resultados devem ser analisados com cautela, pois o material disponível aponta limitações como elevada rotatividade das carteiras, custos de transação capazes de reduzir parte dos lucros, concentração dos testes em mercados específicos e uso predominante da variância como métrica de risco, com menor atenção a medidas como VaR e CVaR \cite{chen2021mean,ma2021portfolio,wang2020portfolio,padhi2022intelligent}.

\section{Estudos sobre média-variância e Markowitz em criptoativos}
\label{sec_estudos_sobre_media_variancia_e_markowitz_em_criptoativos}

Enquanto os estudos anteriores mostram o uso de previsões algorítmicas na otimização de carteiras, a literatura sobre criptoativos permite observar como o modelo média-variância se comporta em um mercado mais volátil. A Teoria Moderna de Portfólio tem sido usada para avaliar os efeitos da diversificação sobre a relação entre risco e retorno em carteiras formadas por criptomoedas ou por combinações entre ativos digitais e ativos tradicionais \cite{brauneis2019cryptocurrency,jeleskovic2024cryptocurrency}. Nesse contexto, o modelo de Markowitz aparece tanto como instrumento para avaliar benefícios de diversificação quanto como método para definir pesos otimizados segundo o critério média-variância \cite{brauneis2019cryptocurrency,jeleskovic2024cryptocurrency}.

A aplicação desse modelo ao mercado de criptoativos exige cautela. Os estudos analisados destacam que esse mercado apresenta elevada volatilidade, retornos não normais, correlações instáveis, dificuldade de previsão da dependência entre ativos e sensibilidade a fatores globais. Essas características podem limitar a aderência das premissas clássicas da média-variância \cite{brauneis2019cryptocurrency,jeleskovic2024cryptocurrency,zouaoui2025portfolio}. Ainda assim, os resultados sobre diversificação sugerem que criptoativos podem contribuir para a redução de risco quando combinados entre si ou quando inseridos em carteiras com ativos tradicionais. \citeonline{brauneis2019cryptocurrency} apontam que a combinação de várias criptomoedas apresenta potencial de redução de risco, ampliando a análise para além da inclusão isolada de uma única moeda digital. \citeonline{jeleskovic2024cryptocurrency}, de modo complementar, indicam que portfólios compostos por criptomoedas e ativos tradicionais apresentam maior Índice de Sharpe em avaliação retrospectiva e desempenho mais estável em análise prospectiva.

Apesar desse potencial, os estudos também apontam fragilidades da otimização média-variância clássica em carteiras de criptomoedas. \citeonline{brauneis2019cryptocurrency} observam que a estratégia de pesos iguais, representada pelo portfólio 1/N, supera criptomoedas individuais e mais de 75\% dos portfólios ótimos de média-variância em termos de Índice de Sharpe e retornos equivalentes de certeza. Esse achado sugere que a otimização matemática tradicional nem sempre apresenta desempenho mais favorável do que estratégias simples de diversificação \cite{brauneis2019cryptocurrency}. De forma próxima, \citeonline{zouaoui2025portfolio} apontam que o modelo de Markowitz pode concentrar a carteira em poucas criptomoedas, deixando a maior parte dos ativos fora da alocação, situação associada, segundo os autores, a limitações no uso do coeficiente de correlação e à geração restrita de portfólios possíveis.

A relação entre risco e retorno reforça essas limitações. \citeonline{zouaoui2025portfolio} reportam que uma carteira baseada no modelo retorno-risco de Markowitz apresentou retorno estimado de 31,15\% e risco de 39,05\%, além de Índice de Sharpe médio pouco atrativo no portfólio selecionado. Os autores associam parte desse risco às fortes oscilações dos preços das criptomoedas e à influência de condições globais sobre o mercado. \citeonline{jeleskovic2024cryptocurrency}, por sua vez, destacam que a imaturidade estrutural do mercado e a dificuldade de prever a estrutura de correlação entre criptomoedas podem gerar underfitting nos modelos de risco, prejudicando a identificação adequada das dependências entre os ativos. Esses achados indicam que a volatilidade e a instabilidade das correlações podem limitar a aplicação isolada da média-variância no mercado cripto \cite{jeleskovic2024cryptocurrency,zouaoui2025portfolio}.

Essas limitações estão relacionadas às próprias premissas do modelo. \citeonline{brauneis2019cryptocurrency} reconhecem a existência de alternativas à otimização média-variância quando os retornos não seguem distribuição normal. \citeonline{zouaoui2025portfolio} também ressaltam que a Teoria Moderna de Portfólio, embora ofereça uma estrutura quantitativa para diversificação e análise da relação entre risco e retorno, pode ter sua aplicação prática restringida em mercados voláteis, não normais e sujeitos a correlações variáveis. Dessa forma, a literatura analisada sugere que Markowitz permanece como referência metodológica de alocação, mas apresenta limitações importantes quando aplicado diretamente a carteiras de criptoativos \cite{brauneis2019cryptocurrency,zouaoui2025portfolio}.

Nesse cenário, a incorporação de técnicas preditivas aparece como uma alternativa para lidar com parte dessas fragilidades. \citeonline{zouaoui2025portfolio} indicam que redes neurais, como LSTM, quando associadas à Teoria Moderna de Portfólio no mercado de criptomoedas, podem reduzir algumas limitações observadas no modelo clássico e gerar Índices de Sharpe mais elevados. Essa discussão reforça a relevância de investigar abordagens híbridas que combinem média-variância e machine learning, uma vez que os estudos apontam tanto o potencial da diversificação em criptoativos quanto os limites da aplicação isolada de Markowitz nesse mercado \cite{brauneis2019cryptocurrency,jeleskovic2024cryptocurrency,zouaoui2025portfolio}.

\section{Estudos híbridos entre Markowitz, machine learning, deep learning e reinforcement learning}
\label{sec_estudos_hibridos_entre_markowitz_machine_learning_deep_learning_e_reinforcement_learning}

A partir das limitações da média-variância aplicada de forma isolada, parte da literatura passou a investigar modelos híbridos que combinam Markowitz com técnicas preditivas e adaptativas. Nos estudos analisados, a Teoria Moderna de Portfólio não é abandonada, mas utilizada como base de alocação à qual são incorporados algoritmos voltados a melhorar estimativas, selecionar ativos ou ajustar dinamicamente os pesos do portfólio \cite{chaweewanchon2022markowitz,du2022mean,lopez2025enhancing}. Essa lógica se aproxima do presente anteprojeto, pois articula otimização média-variância e métodos computacionais voltados à previsão e à tomada de decisão em carteiras financeiras \cite{chaweewanchon2022markowitz,slusarczyk2025optimal,lopez2025enhancing}.

No primeiro grupo de estudos, a combinação entre Markowitz, machine learning e deep learning aparece associada à pré-seleção de ativos e à melhoria dos dados usados na otimização. \citeonline{chaweewanchon2022markowitz} combinam CNN, BiLSTM e modelo média-variância, utilizando previsões para selecionar previamente ações antes da formação do portfólio. \citeonline{du2022mean}, por sua vez, parte da limitação de que alocações média-variância são sensíveis a pequenos erros nas estimativas de médias e covariâncias, propondo o uso de redes A-LSTM para prever portfólios cointegrados antes da otimização. Em ambos os casos, os algoritmos preditivos funcionam como etapa anterior ou complementar ao modelo de Markowitz, buscando melhorar os dados usados na construção da carteira \cite{chaweewanchon2022markowitz,du2022mean}.

A previsão de retornos dentro da estrutura média-variância também aparece como eixo relevante. \citeonline{slusarczyk2025optimal} comparam carteiras construídas com retornos históricos e carteiras baseadas em retornos previstos por ARIMA-GARCH e XGBoost, considerando problemas de otimização voltados ao máximo índice de informação e à mínima variância global. Os achados indicam que o uso de retornos previstos pode melhorar a seleção de portfólio no modelo de Markowitz, mas essa solução não é universal e depende do controle cuidadoso de parâmetros e hiperparâmetros \cite{slusarczyk2025optimal}. De forma complementar, \citeonline{lopez2025enhancing} discutem a integração de técnicas de machine learning à seleção de portfólio de Markowitz, destacando sua aplicação em processos como geração de alfa, gestão de risco e otimização de métricas como o CVaR.

Uma vertente mais dinâmica é observada nos estudos que conectam Teoria Moderna de Portfólio e deep reinforcement learning. \citeonline{jiang2024deep} desenvolvem uma estrutura de DRL para seleção de portfólios em mercados financeiros complexos e de grande dimensão, incorporando aversão ao risco e custos de transação em uma função de recompensa baseada em uma extensão da média-variância de Markowitz. Nessa abordagem, o agente de aprendizado utiliza estados de mercado para treinar o modelo e ajustar a decisão de alocação, considerando risco e custos durante o processo de aprendizagem \cite{jiang2024deep}. \citeonline{jang2023deep} também conectam reinforcement learning e Teoria Moderna de Portfólio, utilizando decomposição tensorial para combinar entradas de diferentes dimensões e permitir alterações dinâmicas nos pesos da carteira conforme tendências de mercado.

Comparativamente, os estudos baseados em machine learning e deep learning concentram-se na melhoria das previsões e dos dados de entrada usados pela média-variância. Já os estudos baseados em reinforcement learning ampliam a discussão para uma lógica de decisão sequencial e de alocação dinâmica da carteira \cite{chaweewanchon2022markowitz,du2022mean,jiang2024deep,jang2023deep}. Em termos de resultados, os modelos preditivos indicam potencial de aprimoramento da média-variância, com efeitos reportados sobre Índice de Sharpe, retorno médio, risco e lucratividade em diferentes mercados acionários \cite{chaweewanchon2022markowitz,du2022mean}. Os modelos de deep reinforcement learning, por sua vez, apresentam resultados mais favoráveis e estáveis em relação a estratégias de referência, além de realizarem alocação dinâmica conforme estados ou tendências de mercado \cite{jiang2024deep,jang2023deep}.

Apesar desses achados, as abordagens híbridas também apresentam limitações. \citeonline{slusarczyk2025optimal} ressaltam que retornos previstos não devem ser tratados como solução universal para a estrutura de Markowitz, pois a eficácia do método depende do controle dos parâmetros e hiperparâmetros adotados. Assim, os resultados sugerem potencial de melhoria na alocação de carteiras por meio de técnicas preditivas e adaptativas, mas também mostram a necessidade de cuidado na configuração dos modelos \cite{slusarczyk2025optimal}.

Em síntese, os estudos analisados evidenciam que a integração entre Markowitz, machine learning, deep learning e reinforcement learning ocorre de diferentes formas: pré-seleção de ativos, previsão de retornos, atualização dos dados usados na otimização e alocação dinâmica por agentes de aprendizado \cite{chaweewanchon2022markowitz,du2022mean,slusarczyk2025optimal,jiang2024deep,jang2023deep,lopez2025enhancing}. Essa diversidade metodológica oferece base para investigar modelos híbridos também no contexto de carteiras de criptoativos, especialmente pela combinação entre uma estrutura clássica de diversificação, baseada em média-variância, e técnicas de aprendizado voltadas à previsão, gestão de risco e ajuste dinâmico da alocação \cite{jiang2024deep,lopez2025enhancing,slusarczyk2025optimal}.

\section{Síntese crítica dos estudos relacionados}
\label{sec_sintese_critica_dos_estudos_relacionados}

Consideradas em conjunto, as subseções anteriores indicam uma tendência de integração entre modelos clássicos de otimização de portfólios e técnicas preditivas baseadas em machine learning, deep learning e reinforcement learning. Essa integração assume diferentes formas, como previsão de preços ou retornos, seleção prévia de ativos, melhoria dos dados usados pela média-variância e ajuste dinâmico da alocação da carteira \cite{chen2021mean,ma2021portfolio,wang2020portfolio,chaweewanchon2022markowitz,du2022mean,slusarczyk2025optimal,jiang2024deep,jang2023deep,lopez2025enhancing}. Nos estudos analisados, a inteligência artificial aparece menos como substituta da Teoria Moderna de Portfólio e mais como recurso complementar para lidar com dados financeiros complexos, não lineares, volumosos e sujeitos a mudanças ao longo do tempo \cite{sutiene2024enhancing,giantsidi2025deep,lopez2025enhancing}.

Em parte significativa dos estudos, os resultados apontam melhorias em retorno, risco, retorno ajustado ao risco ou Índice de Sharpe quando modelos preditivos são incorporados à formação de carteiras. Pesquisas com Random Forest, XGBoost, LSTM, CNN-BiLSTM e A-LSTM indicam que a previsão de retornos ou a pré-seleção de ativos pode produzir resultados mais favoráveis do que estratégias baseadas apenas em dados históricos, sem etapa preditiva ou com alocação tradicional, conforme o desenho de cada estudo \cite{chen2021mean,ma2021portfolio,wang2020portfolio,chaweewanchon2022markowitz,du2022mean,slusarczyk2025optimal}. Em outra direção, os modelos de deep reinforcement learning ampliam essa lógica ao incorporar estados de mercado, aversão ao risco e custos de transação ao processo de decisão, permitindo formas mais dinâmicas de alocação em comparação a estratégias estáticas de média-variância ou Max-Sharpe \cite{jiang2024deep,jang2023deep}.

Contudo, os ganhos não são universais. \citeonline{slusarczyk2025optimal} indicam que o uso de retornos previstos pode aprimorar a seleção de portfólios no modelo de Markowitz, mas depende do controle cuidadoso de parâmetros e hiperparâmetros. No mercado de criptomoedas, \citeonline{brauneis2019cryptocurrency} apontam que a estratégia 1/N superou criptomoedas individuais e parcela expressiva dos portfólios ótimos de média-variância em termos de Índice de Sharpe e retornos equivalentes de certeza. Esses achados sugerem que a sofisticação metodológica não assegura, isoladamente, melhor desempenho em todos os contextos, sobretudo no mercado de criptoativos, marcado por alta volatilidade, retornos não normais e correlações instáveis \cite{brauneis2019cryptocurrency,jeleskovic2024cryptocurrency,zouaoui2025portfolio}.

As limitações metodológicas reforçam essa cautela. Nos estudos sobre otimização com previsão, o material analisado aponta restrições associadas à rotatividade das carteiras e aos custos de transação, fatores que podem reduzir parte dos ganhos associados às estratégias preditivas \cite{chen2021mean,ma2021portfolio,wang2020portfolio,padhi2022intelligent}. Revisões recentes também destacam que modelos híbridos e de deep learning podem exigir maior custo computacional e planejamento metodológico cuidadoso \cite{giantsidi2025deep}. Diferenças nos dados, no pré-processamento e nas estratégias de validação dificultam a comparação entre resultados e aumentam o risco de overfitting \cite{giantsidi2025deep}. De modo complementar, \citeonline{lopez2025enhancing} apontam que o overfitting permanece como desafio relevante em estratégias de portfólio baseadas em machine learning, afetando a capacidade dos modelos de manter desempenho robusto em diferentes condições de mercado. Essa dificuldade de generalização também afeta a média-variância, pois carteiras otimizadas com bom desempenho histórico podem apresentar desempenho inferior em condições futuras de mercado \cite{lopez2025enhancing}.

Outro limite relevante é a interpretabilidade. Embora a inteligência artificial amplie a capacidade de processar grandes volumes de dados financeiros e identificar padrões complexos, parte dos modelos de machine learning opera como ``caixa-preta'', dificultando a compreensão das razões associadas às decisões de alocação \cite{sutiene2024enhancing}. \citeonline{giantsidi2025deep} também apontam que métodos de explicabilidade, como abordagens de XAI, ainda precisam ser associados a validações específicas do domínio financeiro. Assim, a literatura sugere que o avanço dos modelos aplicados a portfólios deve considerar não apenas arquiteturas mais complexas, mas também sistemas interpretáveis, adaptativos e compatíveis com as restrições práticas dos mercados reais \cite{sutiene2024enhancing,giantsidi2025deep}.

No caso dos criptoativos, a comparação entre os estudos evidencia uma lacuna específica. A literatura sobre média-variância em criptomoedas aponta potencial de diversificação, principalmente quando diferentes criptoativos são combinados entre si ou com ativos tradicionais \cite{brauneis2019cryptocurrency,jeleskovic2024cryptocurrency}. Ao mesmo tempo, os estudos indicam limitações importantes da aplicação direta de Markowitz nesse mercado, em razão da volatilidade elevada, da não normalidade dos retornos, da instabilidade das correlações e da dificuldade de estimar adequadamente a estrutura de dependência entre os ativos \cite{brauneis2019cryptocurrency,jeleskovic2024cryptocurrency,zouaoui2025portfolio}. Como parte relevante dos estudos híbridos analisados se concentra em ações, índices ou mercados financeiros tradicionais, permanece espaço para investigar de forma mais específica modelos híbridos aplicados à diversificação de portfólios em criptoativos \cite{chaweewanchon2022markowitz,du2022mean,slusarczyk2025optimal,jiang2024deep,jang2023deep}.

Dessa forma, o presente anteprojeto se insere nessa lacuna ao examinar a combinação entre previsão por machine learning e otimização média-variância na construção de carteiras de criptoativos. Essa investigação deve considerar as limitações associadas à instabilidade do mercado cripto e à aplicação direta de Markowitz nesse contexto \cite{brauneis2019cryptocurrency,jeleskovic2024cryptocurrency,zouaoui2025portfolio}, aos custos de transação e à rotatividade da carteira em estratégias preditivas \cite{chen2021mean,ma2021portfolio,wang2020portfolio,padhi2022intelligent}, à sensibilidade a parâmetros e hiperparâmetros \cite{slusarczyk2025optimal}, ao risco de overfitting e à dificuldade de generalização \cite{giantsidi2025deep,lopez2025enhancing}, bem como à baixa interpretabilidade de parte dos modelos de machine learning \cite{sutiene2024enhancing,giantsidi2025deep}. Assim, a proposta mantém o foco na integração entre previsão algorítmica e diversificação matemática, direcionando essa discussão ao contexto específico dos criptoativos e à avaliação crítica de sua viabilidade em estratégias híbridas de portfólio.
